\begin{afterword}
\summaryhead

The post starts from the point that in real life people must search for their options and their evidence, and must decide when to stop, as when house-hunting. Gilovich had described how people ask whether the evidence allows a welcome conclusion but whether it compels an unwelcome one. By analogy, the author suggests two biases in search. With a hidden motive to accept the current best option, we stop too soon: motivated stopping. With a hidden motive to reject it, we keep looking and ask for more evidence: motivated continuation.

The example is R. A. Fisher, who insisted that smoking had not been shown to cause lung cancer, while working as a consultant to tobacco firms and smoking himself. Asking for more evidence can look virtuous, the author says, but evidence is costly and slow, and one can always change one's mind later. The post ends with a test: suspect motivated stopping when cheap evidence is left ungathered, and motivated continuation when the evidence demanded is expensive and slow.

\responsehead

The distinction is useful, and the closing test is practical. ``The Bottom Line'' and ``Rationalization'' warned against deciding first and gave no way to catch oneself doing it. This post gives a sign one can observe: the cost and speed of the evidence one skips, or demands, before an uncomfortable decision.

The idea is offered as new (``I suggest''), and it was not. The stopping half is in the author's own ``The Third Alternative,'' five months earlier, with the same house example. Both halves were in the psychology literature. Kruglanski had argued, in Ditto and Lopez's summary, that motives can hasten or delay the ``freezing'' of a search. Ditto and Lopez, in a 1992 paper titled ``Motivated Skepticism,'' measured both: people needed less information to reach a conclusion they wanted, and people given an unwelcome medical result took longer to accept it and were more likely to retest. The post uses the term ``motivated skepticism,'' the title of that paper, and cites none of the work that had tested its proposal.

The Fisher example is sound in its main facts and doubtful in one detail. Fisher did argue that the association did not show causation, was a paid consultant to the tobacco industry, and smoked. The motive is offered with ``maybe,'' and a historian has suggested the same. But I could not confirm that Fisher ``testified to Congress.'' Fisher's arguments appeared in print, in letters and articles collected in a pamphlet. Fisher's hypothesis did reach Congress, in 1965 and 1969, through an industry consultant, after Fisher's death.

The footnote that contrasts Fisher the frequentist with ``more reasonable'' Bayesians cites Judea Pearl, whose own view is different. Pearl holds that causal claims cannot be drawn from associations alone by any school of probability; on that view Fisher's slogan was right, and his error, in Stolley's account, lay in not examining all the data available.

Finally, the two halves are unevenly supported. Continuation has the post's one example. Stopping gets a list of places where it ``appears'' and no case.

In short: a useful distinction with a practical test, presented as a new suggestion although it had been proposed and measured, and illustrated by a case whose main facts are correct.
\end{afterword}
